\subsection{Absolute Tate cohomology: conventions and the missing $2$-torsion}
\label{sec:real-weil-E4}

\paragraph{E.4a (Proved): all twists over a finite field.}
Let $q=p^a$. Use Lichtenbaum--Geisser Weil-\'etale cohomology of
$\operatorname{Spec}\mathbb F_q$, with the ordinary \'etale motivic complex
$\mathbb Z(n)$ for $n\geq0$ and the explicit convention
\[
 \mathbb Z(n)=(\mathbb Q/\mathbb Z)'(n)[-1]\quad(n<0),
 \qquad(\mathbb Q/\mathbb Z)'=\bigoplus_{\ell\ne p}\mathbb Q_\ell/\mathbb Z_\ell.
 \tag{E.1}
\]
The arithmetic Frobenius generator acts on roots of unity by their $q$th
power. With this convention, for every $i,n\in\mathbb Z$,
\[
 H_W^i(\mathbb F_q,\mathbb Z(n))\simeq
 \begin{cases}
  \mathbb Z,& n=0,\ i=0,1,\\
  \mathbb Z/(q^{|n|}-1),& n\ne0,\ i=1,\\
  0,&\text{otherwise}.
 \end{cases}
 \tag{E.2}
\]
The finite cyclic-group isomorphism is not a specified choice of generator.
For $n\ne0$ the $p$-primary torsion is zero; for $n=0$ both nonzero
groups are free and their $p$-adic completions are $\mathbb Z_p$.

\begin{proof}
For $n=0$, the Weil group is $\mathbb Z$ acting trivially on $\mathbb Z$;
its two-term group-cohomology resolution gives $\mathbb Z$ in degrees
$0,1$ and zero elsewhere. For the other twists we justify both finite
generation and the comparison used in the computation.
\citet[Proposition 9.4]{D_Geisser2004} gives the integral comparison
$K(\mathbb P^1,n)$ for every $n$: for a product of one curve its condition
$n\leq1$ or $n\geq d-1$ covers all integers. Theorem 1.2 in that paper,
applied also to the complementary twist, gives finite generation.
Pullback along $\mathbb P^1\to\operatorname{Spec}\mathbb F_q$ and along a
rational point compose to the identity. Thus the point is a direct
summand, and inherits both assertions, compatibly with the comparison
maps. This avoids presuming a theorem for an empty product of curves.

For $\ell\ne p$ continuous cohomology is computed by
$[\mathbb Z_\ell(n)\xrightarrow{\,1-q^n\,}\mathbb Z_\ell(n)]$ in degrees
$0,1$. The differential is nonzero if $n\ne0$, giving
$\mathbb Z_\ell/(1-q^n)$ in degree one and no other group.
For positive $n$ the logarithmic de Rham--Witt sheaf
$W_r\Omega^n_{\mathbb F_q,\log}$ is zero, so the $p$-adic comparison
has zero groups. For negative $n$ this follows directly from (E.1).
A finitely generated abelian group is determined by these completions:
there is no free part, each $\ell$-primary part is cyclic of order
$\ell^{v_\ell(q^{|n|}-1)}$, and the $p$-part is zero. Their direct sum
is the cyclic group in (E.2). Equivalently, using (E.1) and the rational
vanishing, one computes the invariants and coinvariants of the divisible
group $(\mathbb Q/\mathbb Z)'(n)$: multiplication by $q^n-1$ is
surjective and its kernel has the indicated order.
\end{proof}

The negative-twist convention is exactly the one printed at the beginning
of \S7 of \citet{D_Geisser2004}; Theorem 7.1 supplies its comparison with
\'etale cohomology. For $n<0$ the literal expression
$\mathbb Z/(1-q^n)\mathbb Z$ on 15:173 is not an integral quotient as
written. Its consistent interpretation is
$\mathbb Z[1/p]/(1-q^n)\mathbb Z[1/p]$, canonically isomorphic to the
cyclic group of order $q^{|n|}-1$. This is a correction of notation,
not an assertion that negative-twist motivic cohomology is independent
of conventions.

\paragraph{E.4b (Refuted, with corrected groups): exact agreement with the integral number-ring table.}
The values on 15:173 do not agree integrally with the standard Weil-\'etale
normalization. Let $\overline X$ denote the Artin--Verdier compactification
of $X=\operatorname{Spec}\mathbb Z$, and use the derived Weil-\'etale
complex of \citet[Definition 3.6]{D_FlachMorin2018}. The groups are
\[
\begin{array}{c|cccc|c}
 &H_W^0&H_W^1&H_W^2&H_W^3&H_W^i\ (i\notin\{0,1,2,3\})\\ \hline
 \mathbb Z(0)&\mathbb Z&0&0&\mathbb Z/2&0\\
 \mathbb Z(1)&0&\mathbb Z/2&0&\mathbb Z&0.
\end{array}
 \tag{E.3}
\]
Thus $H_W^3(\overline X,\mathbb Z(1))=\mathbb Z$ is the compactified
Weil-\'etale trace group. It is not ordinary motivic cohomology with
$\mathbb Z(1)=\mathbb G_m[-1]$ on the uncompactified \'etale site.

\begin{proof}
The number-ring computation of
\citet[\S5.8.3, pp.~1503--1504]{D_FlachMorin2018} gives, in twist $1$,
$\mathcal O_F^\times$, $\operatorname{Cl}(\mathcal O_F)$, and $\mathbb Z$
in degrees $1,2,3$. For $F=\mathbb Q$ these are
$\{\pm1\}$, $0$, and $\mathbb Z$. In twist $0$, the same table gives
$\mathbb Z$, $0$,
$\operatorname{Hom}(\mathcal O_F^\times,\mathbb Z)
\oplus\operatorname{Cl}(\mathcal O_F)^D$, and
$(\mathcal O_F^\times)_{\rm tors}^D$ in degrees $0,1,2,3$.
Substituting $\mathbb Z^\times=\{\pm1\}$ proves (E.3).
The trace is also the specialization $d=1,n=1$ of their Theorem 3.22.
The hypotheses of their construction are known for number rings, as
stated at the start of \S5.8.3, and are not being assumed for arbitrary
arithmetic schemes here.
\end{proof}

The contradiction to the asserted exact integral agreement is already
$H_W^1(\overline X,\mathbb Z(1))=\mathbb Z/2\ne0$. Likewise the
unqualified twist-zero assertion misses a group of order two in degree
three. These groups disappear after inverting $2$ and rationally; the
rational $0,0,0,\mathbb Q$ table in R15 survives this check.
The source attached in the archive, \citet[Theorem 5.10]{D_Lichtenbaum2005},
already gives $H^3(\overline X_W,\mathbb Z)=\mu(\mathbb Q)^D=\mathbb Z/2$.
Its theory with the constant sheaf only in low degrees should not be
silently equated with a later bounded complex in every degree. The
published number-ring paper \citep{D_Lichtenbaum2009} treats the
constant sheaf and compact support at infinity, not a construction of
all the twists requested here.

\paragraph{E.4c (Proved for the localization extension constructed below): the displayed open-subset calculation.}
Let $S$ be a finite nonempty set of rational primes and retain the real
place when removing $S$ from $\overline X$. If the open complex is defined
by the localization triangle with purity
$R i_p^!\mathbb Z(1)=\mathbb Z(0)[-2]$ and trace sending each finite-field
generator to $1$, its twist-one groups are
\[
\begin{split}
 H_W^1(\overline X-S,\mathbb Z(1))&\simeq
                \mathbb Z/2\oplus\mathbb Z^S,\\
 H_W^2(\overline X-S,\mathbb Z(1))&\simeq
                \ker(\mathbb Z^S\xrightarrow{\sum}\mathbb Z),\\
 H_W^i(\overline X-S,\mathbb Z(1))&=0\quad(i\ne1,2).
\end{split}
 \tag{E.4}
\]
Indeed the first piece of the localization sequence is
$0\to\mathbb Z/2\to H_W^1\to\mathbb Z^S\to0$, which splits as a
sequence of abelian groups; the next is
$0\to H_W^2\to\mathbb Z^S\to\mathbb Z\to H_W^3\to0$.
The sum map is surjective. This proves (E.4) for exactly
the localization normalization stated and constructed in E.4e, and proves the rational calculation
on 15:174--178. In any extension satisfying the corresponding
purity/localization triangle in twist $n\ne1$, its residue terms vanish
rationally by (E.2), and the restriction invariance written on 15:174 follows.
The margin's $n\ne0$ on 15:173 is therefore not the correct condition.

E.4o below identifies this bounded localization extension with an
explicit rational correction to the Artin--Verdier etale motivic complex,
including all finite coefficients; E.4p extends it to every twist on
this family of opens. A comparison with every integral Weil-\'etale
construction and every real-place convention is not asserted. The assertion ``a cohomology theory with exactly his values''
is refuted under the standard normalization by (E.3); without axioms
for a cohomology theory it is not a uniqueness statement. Lead E2 is
correct as a description of the finite-field and rational trace pattern,
but false as an assertion of exact integral identification with the
attached 2005 theory.
An additional modification at $2$ is not ruled out by this discrepancy;
E.4f constructs one on these particular complexes and checks its
localization and proper duality, while spelling out its lost comparison.
The phrase \textit{petite modification $l=2$ en dim $0$} on 15:178 occurs
beside a Kummer computation, after the same page explicitly retained
$\mathbb Z/2$ in $H^0(\overline U,\mathbb G_m)$. It is therefore not a
definition that permits us to erase the sign group from every complex.

\paragraph{E.4d (Refuted under the coefficient comparison already used in the manuscript): the missing sign cannot be ignored.}
There is an elementary obstruction to retaining the exact integral
twist-one table while also retaining the usual coefficient triangle and
degree-zero Kummer comparison. Let $C=R\Gamma(\overline X,\mathbb Z(1))$
be a complex of abelian groups. Suppose
$H^0(C)=H^1(C)=0$, as on 15:173, and suppose reduction modulo $2$ is
computed by the distinguished triangle
\[
 C\xrightarrow{2}C\longrightarrow
       R\Gamma(\overline X,\mathbb Z/2(1))\longrightarrow C[1].
\]
Its long exact sequence forces
$H^0(\overline X,\mathbb Z/2(1))=0$.
The degree-zero Kummer identification instead gives
$H^0(\overline X,\mathbb Z/2(1))=\mu_2(\mathbb Z)=\{\pm1\}$.
These assumptions are inconsistent. The same obstruction appears on
$\operatorname{Spec}\mathbb Z[1/2]$, where the ordinary \'etale Kummer
sequence is exact, if one omits the sign torsion from its twist-one
$H^1$.

The standard corrected group $H^1(C)=\mathbb Z/2$ accounts for the
missing finite-coefficient degree-zero group through the connecting
homomorphism. Thus any further real-place modification giving the
literal table must explicitly change this Kummer comparison or its
coefficient triangle. It cannot simultaneously preserve all the
comparisons written on 15:173--178. This is a refutation with stated
axioms, and does not assume that every imaginable assignment of groups
is a cohomology theory with those axioms.

\paragraph{E.4e (Proved): the normalized Gysin map and an explicit open complex.}
There is a canonical bounded derived localization extension of the
verified compactified twist-one complex to the opens
$\overline U_S=\overline{\operatorname{Spec}\mathbb Z}-S$.
It has the groups (E.4), and its construction does not require applying
a flat-pullback theorem to the nonflat immersion
$\operatorname{Spec}\mathbb F_p\hookrightarrow\operatorname{Spec}\mathbb Z$.

\begin{proof}
Put $C_X^n=R\Gamma_W(\overline X,\mathbb Z(n))$ and
$C_p^0=R\Gamma_W(\mathbb F_p,\mathbb Z)$. The already verified groups
and the vanishing of $\operatorname{Ext}_{\mathbb Z}^j$ for $j>1$ give
canonical decompositions in $D(\mathbb Z)$
\[
 \begin{aligned}
 C_X^0&=\mathbb Z\oplus(\mathbb Z/2)[-3],&
 C_X^1&=(\mathbb Z/2)[-1]\oplus\mathbb Z[-3],\\
 C_p^0&=\mathbb Z\oplus\mathbb Z[-1].
 \end{aligned}
 \tag{E.4.1}
\]
For the last splitting the only possible ambiguity is
$\operatorname{Ext}^1(\mathbb Z,\mathbb Z)=0$; in the first two
the degrees are separated farther, so the corresponding ambiguity
groups also vanish. The identifications on $H^0$ are given by the
constant section $1$, and those on the trace groups by the positive
arithmetic Frobenius generator and the compactified trace.

There is exactly one map $\rho_p:C_X^0\to C_p^0$ inducing the identity
on $H^0$: all other components vanish by (E.4.1).
The perfect duality of \citet[Theorem 3.22]{D_FlachMorin2018} identifies
$C_X^1$ with $R\operatorname{Hom}(C_X^0,\mathbb Z[-3])$.
For the finite field, the two-term resolution in E.4a has the
circle duality
$C_p^0\simeq R\operatorname{Hom}(C_p^0,\mathbb Z[-1])$.
Dualizing $\rho_p$ therefore gives a specified Gysin map
\[
 \gamma_p:C_p^0[-2]\longrightarrow C_X^1.
 \tag{E.4.2}
\]
Under (E.4.1), it is zero on the degree-two copy of $\mathbb Z$
and the identity on the degree-three copy.

Define the extension by the cofiber, using ordinary complexes so that
all inclusions of finite sets give coherent maps:
\[
 C_S^1=\operatorname{Cone}
 \left(\bigoplus_{p\in S}C_p^0[-2]
                 \xrightarrow{\ \sum\gamma_p\ }C_X^1\right).
 \tag{E.4.3}
\]
An explicit representative is
\[
 \begin{gathered}
 (C_S^1)^1=\mathbb Z/2\oplus\mathbb Z^S,\qquad
 (C_S^1)^2=\mathbb Z^S,\qquad (C_S^1)^3=\mathbb Z,\\
 d^1=0,\qquad d^2=\sum ,
 \end{gathered}
 \tag{E.4.4}
\]
with all other terms zero. For $S\subset T$, use the coordinate
inclusions in degrees one and two and the identity in degree three.
These are chain maps and compose strictly. Their cofiber triangle
is the localization triangle for the primes $T-S$, as is seen either
by the cone construction or by its degreewise split quotient.
Thus (E.4.3) realizes the normalization of E.4c, including iterated
localization, rather than merely assuming that its long exact sequence
exists. The degree-one free group has the familiar unit identification
$(e,(a_p))\mapsto(-1)^e\prod_{p\in S}p^{a_p}$.
\end{proof}

For completeness, the real-place correction also affects the finite
coefficient calculation. For $m\geq1$, set $d=\gcd(m,2)$ and take
$S\ne\varnothing$. Derived reduction of (E.4.4) gives
\[
 H^i(C_S^1\otimes^{L}\mathbb Z/m)=
 \begin{cases}
 \mathbb Z/d,&i=0,\\
 \mathbb Z/d\oplus(\mathbb Z/m)^S,&i=1,\\
 \ker((\mathbb Z/m)^S\xrightarrow{\sum}\mathbb Z/m),&i=2,\\
 0,&\text{otherwise}.
 \end{cases}
 \tag{E.4.5}
\]
Indeed resolve the $\mathbb Z/2$ in degree one by
$\mathbb Z\xrightarrow{2}\mathbb Z$ in degrees zero and one;
the other terms are free. The sum map is split surjective when $S$
is nonempty. For $S=\varnothing$ there are the same sign groups in
degrees zero and one, together with $\mathbb Z/m$ in degree three.
Thus for even finite coefficients the sign contributes in two degrees.
When $2$ is invertible on $U_S$, this also exposes the sign term in
the ordinary Kummer calculation of units modulo squares.
The corresponding $2$-adic integral degree-one group retains
$\mathbb Z/2$, whereas the inverse-limit degree-zero group is zero.
These distinctions are more precise than deleting a degree-zero term
without changing its integral source.

The construction just made is an explicitly defined bounded derived
extension of the compactified theory, with the real place retained.
It is not a proof that the original unbounded number-ring Weil topos,
every compact-support convention, and this extension are all equivalent.
The equality of those constructions in all twists would require a
separate comparison theorem.

\paragraph{E.4f (Proved at the stated complex level): an explicit sign-deleted modification.}
At $\overline X$ define
\[
 \begin{aligned}
 C_X^{1,+}&=\operatorname{Cone}((\mathbb Z/2)[-1]\to C_X^1)
                 =\mathbb Z[-3],\\
 C_X^{0,+}&=\tau_{\leq0}C_X^0=\mathbb Z .
 \end{aligned}
 \tag{E.4.6}
\]
The first arrow is the canonical truncation inclusion.
These two complexes are perfectly dual with shift $[-3]$.
Retain $C_p^0$ at each finite point and use the composite of (E.4.2)
with $C_X^1\to C_X^{1,+}$ as its Gysin map. Define
$C_S^{1,+}$ by (E.4.3) with that target. It is (E.4.4) with
the $\mathbb Z/2$ summand deleted. Equivalently,
\[
 C_S^{1,+}=\operatorname{Cone}((\mathbb Z/2)[-1]\to C_S^1).
\]
All restriction maps respect this summand, so the same localization
triangles hold. Consequently (E.4.6) gives exactly the two integral
tables on 15:173, and at nonempty $S$ gives
$H^1=\mathbb Z^S$, $H^2=\ker(\sum:\mathbb Z^S\to\mathbb Z)$
and no other group. The change is entirely $2$-primary.
Under the unit description above, degree one consists of positive
rational $S$-units.

This construction proves that the displayed integral tables can be
realized by a specified modification of the bounded complexes; it
does not evade E.4d. Its reduction modulo $2$ has no sign in degree
zero and hence is not the usual Kummer complex. It is also not,
without further definitions and comparison, a multiplicative motivic
cohomology theory for all schemes and all twists. In particular the
short margin on 15:178 does not uniquely prescribe (E.4.6).
The standard integral identification in Lead E2 remains false, while
the finite-field circle and compactified trace analogies survive with
these precisely stated alternatives.

\paragraph{E.4g (Proved): every rational twist over the compactified number ring.}
The rational answer extends beyond twists zero and one. With the
compactified convention of E.4b, for all integers $i,n$ we have
\[
 H_W^i(\overline{\operatorname{Spec}\mathbb Z},\mathbb Q(n))
 \simeq
 \begin{cases}
 \mathbb Q,&(n,i)=(0,0)\text{ or }(1,3),\\
 \mathbb Q,&n\geq3\text{ odd and }i=1,\\
 \mathbb Q,&n\leq-2\text{ even and }i=2,\\
 0,&\text{otherwise}.
 \end{cases}
 \tag{E.4.7}
\]
We specify the argument, since the negative twists distinguish ordinary
rational motivic cohomology from the duality-completed Weil complex.
For $n>1$, the number-ring computation of
\citet[\S5.8.3, pp.~1503--1504]{D_FlachMorin2018} identifies degree one
with the rationalization of $H^1_{\mathrm{et}}(\operatorname{Spec}\mathbb Z,
\mathbb Z(n))$; the degree-two group is finite and the other displayed
groups are $2$-torsion. The regulator in that same computation identifies
the realification of degree one with
$H^0(\mathbb R,(2\pi i)^{n-1}\mathbb R)$.
Complex conjugation acts here by $(-1)^{n-1}$, so this space has dimension
one precisely for odd $n$. Finite generation then gives the claimed
rational dimensions. The underlying rank theorem is
\citet[Proposition 12.2, p.~271]{D_Borel1974}: for $j\geq2$,
$K_j(\mathbb Z)\otimes\mathbb Q$ has rank one when $j\equiv1\pmod4$
and rank zero otherwise. The motivic degree-one piece has
$j=2n-1$, giving exactly odd $n\geq3$.

The perfect duality of \citet[Theorem 3.22]{D_FlachMorin2018},
after rationalization, gives
\[
 H_W^i(\overline X,\mathbb Q(n))
 \simeq
 \operatorname{Hom}_{\mathbb Q}
 \bigl(H_W^{3-i}(\overline X,\mathbb Q(1-n)),\mathbb Q\bigr).
\]
For $n<0$, the positive twist $1-n>1$ has rank one exactly when
$1-n$ is odd, which is exactly negative even $n$; its degree one
becomes degree two. E.4b supplies $n=0,1$. This proves every case
of (E.4.7). Thus the question whether all other rational groups vanish
has answer no, already in twist three, and also in negative even twists.
The rank theorem used here is the published 1974 theorem. The separately
opened regulator paper \citep{D_Borel1977}, received in November 1976,
was published in 1977; its 1980 erratum \citep{D_Borel1980} concerns
discriminant and complex-sign normalizations, not the rank formula.

There is also a coherent rational extension to all opens with the real
place retained. For $n\ne1$ set
$C_{S,\mathbb Q}^n=C_X^n\otimes\mathbb Q$, with identity restriction
maps; for $n=1$ rationalize (E.4.3). Since
$C_p^{n-1}\otimes\mathbb Q=0$ for $n\ne1$ by E.4a, these assignments
satisfy every finite-prime localization triangle and their iterated
compatibilities. Moreover any extension satisfying those triangles has
the same restriction isomorphisms for $n\ne1$. For $n=1$ and
$S\ne\varnothing$, its only groups are $\mathbb Q^S$ in degree one
and the sum-zero subspace of $\mathbb Q^S$ in degree two.
This settles the entire rational table and the rational restriction
calculation. E.4l--p below construct a compatible integral extension
on this specified family of opens, with all finite-coefficient
comparisons and the necessary wild support. They do not assert a
universal identification among different open Weil-etale topoi or
an arbitrary-scheme multiplicative theory.

\paragraph{E.4h (Proved, with the characteristic-prime restriction): the continuous tables in R15.}
Let $\Gamma=\operatorname{Gal}(\overline{\mathbb F}_q/\mathbb F_q)
\simeq\widehat{\mathbb Z}$, with arithmetic generator $f$, and let $M$
be a finitely generated $\mathbb Z_\ell$-module with continuous
$\Gamma$-action. Continuous cohomology here means the derived inverse
limit of the discrete cohomology of $M/\ell^rM$. Then
\[
 R\Gamma_{\mathrm{cts}}(\Gamma,M)
 \simeq[\,M\xrightarrow{1-f}M\,]
 \quad\text{in degrees }0,1.
 \tag{E.4.8}
\]
Indeed a continuous cocycle into a finite module is determined by its
value at $f$; every value extends continuously after taking a finite
cyclic quotient whose order is divisible by the action order and the
module exponent. Coboundaries give $(f-1)M$.
The usual cyclic resolution, followed by inflation through larger
cyclic quotients, has zero cohomology in degrees at least two: the
periodic classes in degrees at least two are killed after multiplying the
quotient order by a sufficiently large power of $\ell$.
This proves the two-term calculation for each finite quotient.
Their cohomology groups are finite, so the inverse systems satisfy
the Mittag--Leffler condition and have zero $\varprojlim^1$.
Taking the derived inverse limit proves (E.4.8).
For a finite-dimensional continuous $\mathbb Q_\ell$-representation,
compactness of $\Gamma$ gives an invariant lattice. Continuous cochains
have bounded image, so passage from that lattice to its rationalization
commutes with the cochain complex and gives the same formula over
$\mathbb Q_\ell$.

For $\ell\ne p$ we may apply this to the genuine cyclotomic lattice
$M=\mathbb Z_\ell(n)$, where $f$ acts as the unit $q^n$.
If $n=0$, degrees zero and one are both $\mathbb Z_\ell$; if $n\ne0$,
degree zero is zero and degree one is
$\mathbb Z_\ell/(1-q^n)$. Higher degrees vanish.
With $\mathbb Q_\ell(n)$ coefficients, both degrees are
$\mathbb Q_\ell$ for $n=0$ and all groups vanish for $n\ne0$.
These are precisely the continuous tables after stating their topology.

The characteristic prime cannot be inserted into this rank-one
cyclotomic formula. For $n\ne0$, multiplication by $q^n$ is not an
automorphism of $\mathbb Z_p$: it is a nonunit for $n>0$ and is not
even an integral endomorphism for $n<0$. Nor can it give a continuous
one-dimensional $\mathbb Q_p$-representation of $\Gamma$, since the
valuation of a compact subgroup of $\mathbb Q_p^\times$ is zero.
In characteristic $p$ the positive-degree logarithmic de Rham--Witt
term at a finite field is zero, as used in E.4a, and the negative-twist
convention excludes the $p$-part. A crystalline Frobenius with nonunit
eigenvalue belongs to a different coefficient category. Thus R15's
prime-to-$p$ continuous table is confirmed, while an unqualified
extension of that table to a free $p$-adic cyclotomic lattice is false.

\paragraph{E.4i (Proved): the Brauer injectivity and local descent used in R15.}
We also check the separate Brauer-draft assertion rather than importing
it from the earlier reading. Let $X$ be an integral noetherian scheme
whose strict local rings are factorial, let $K$ be its function field,
and let $j:\operatorname{Spec}K\to X$ be the generic-point morphism.
Then
\[
 H^2_{\mathrm{et}}(X,\mathbb G_m)
       \longrightarrow H^2_{\mathrm{et}}(K,\mathbb G_m)
       =\operatorname{Br}(K)
 \quad\text{is injective}.
 \tag{E.4.9}
\]
In particular $\operatorname{Br}(X)\to\operatorname{Br}(K)$ is
injective for the Azumaya Brauer group. These are
\citet[Corollaries 1.8 and 1.10, p.~293]{D_Grothendieck1966BrauerII}
in the opened Bourbaki version; its numbering is not the draft's
Proposition 5.4 and Corollary 5.5.

Here is the proof in the stated setting. On $X_{\mathrm{et}}$,
factoriality of strict local rings gives the exact divisor sequence
\[
 0\longrightarrow\mathbb G_m\longrightarrow j_*\mathbb G_{m,K}
   \longrightarrow
   \bigoplus_{x\in X^{(1)}}(i_x)_*\mathbb Z\longrightarrow0.
\]
The last description is the equality of Cartier and Weil divisors
on strict localizations. For each $x$ the Leray edge map embeds
$H^1(X,(i_x)_*\mathbb Z)$ into
$H^1(k(x),\mathbb Z)=
\operatorname{Hom}_{\mathrm{cts}}(\Gamma_{k(x)},\mathbb Z)=0$;
compactness and torsion-freeness prove the last equality.
Cohomology on the noetherian small \'etale site commutes with the
filtered union of finite subsums. Thus the divisor sheaf has zero
$H^1$, and the displayed sequence injects $H^2(X,\mathbb G_m)$ into
$H^2(X,j_*\mathbb G_{m,K})$.
Moreover $R^1j_*\mathbb G_{m,K}=0$: its stalks are limits of Picard
groups of finite products of separable field extensions of $K$, which
vanish by Hilbert 90. The Leray sequence now injects the latter $H^2$
into $H^2(K,\mathbb G_m)$, proving (E.4.9).
The usual boundary from the exact sequence
$1\to\mathbb G_m\to\mathrm{GL}_r\to\mathrm{PGL}_r\to1$
embeds Azumaya Brauer classes into $H^2(X,\mathbb G_m)$: a zero
boundary means the algebra is an endomorphism algebra of a vector
bundle and hence Brauer-trivial. This proves the Azumaya assertion.

For a local ring $R$, a class
$\xi\in H^2_{\mathrm{et}}(\operatorname{Spec}R,\mathbb G_m)$
comes from an Azumaya algebra if and only if it vanishes after a finite
\'etale faithfully flat extension. This is
\citet[Proposition 2.4, p.~298]{D_Grothendieck1966BrauerII}.
For clarity, the descent direction has an explicit construction.
Refine such an extension to a finite \'etale Galois torsor $B/R$ with
group $G$. The ring $B$ is semilocal, so $\operatorname{Pic}(B)=0$.
The Hochschild--Serre sequence therefore represents $\xi$ by a cocycle
$c\in Z^2(G,B^\times)$. The crossed product
\[
 A_c=\bigoplus_{g\in G}B u_g,
 \qquad (b u_g)(b'u_h)=b\,g(b')c(g,h)u_{gh},
\]
is an $R$-algebra; the cocycle equation is exactly associativity.
After the faithfully flat base change to $B$, the regular crossed
product representation identifies it with a matrix algebra of degree
$|G|$. Hence it is Azumaya, and its descent cocycle gives $\xi$
(replace $c$ by its inverse if the opposite boundary convention is
chosen). Conversely an Azumaya algebra over a local ring has a finite
\'etale maximal subalgebra which splits it: choose a separable maximal
subalgebra in the residue central simple algebra, lift a generator
with separable reduced characteristic polynomial, and use its unit
discriminant. This proves the criterion. It does not assert that every
arbitrary \'etale splitting cover admits a finite \'etale refinement;
the finite condition is the actual hypothesis in this equivalence.

\paragraph{E.4j (Refuted): integral all-twist purity with the verified finite-field groups and ordinary finite coefficients.}
There is a further obstruction at a finite residue prime, independent
of the omitted real sign in E.4b. The following requirements
cannot hold simultaneously. Let
\[
 \overline X=\overline{\operatorname{Spec}\mathbb Z},\qquad
 \overline U=\overline X-\{2\},\qquad z=\operatorname{Spec}\mathbb F_2,
\]
where the real place is retained, and suppose complexes $C_T^n$ have:
the finite-field values E.4a at $z$; the proper twist-zero complex
$C_{\overline X}^0$ of E.4b (or its sign-deleted version E.4f);
an integral purity/localization triangle
\[
 C_z^{-1}[-2]\longrightarrow C_{\overline X}^0
       \longrightarrow C_{\overline U}^0\longrightarrow C_z^{-1}[-1];
 \tag{E.4.10}
\]
and a finite-coefficient comparison
\[
 C_{\overline U}^0\otimes^L\mathbb Z/2
   \simeq R\Gamma(\overline U_{\mathrm{et}},\mathbb Z/2),
 \tag{E.4.11}
\]
with the constant sheaf on the Artin--Verdier site. Already their
degree-one consequences contradict one another.

\begin{proof}
For $q=2,n=-1$, E.4a says that every cohomology group of $C_z^{-1}$
is zero: its possible degree-one group has order $2^1-1=1$.
Thus (E.4.10) makes restriction a quasi-isomorphism.
But E.4b gives
\[
 C_{\overline X}^0\simeq
 \mathbb Z\oplus(\mathbb Z/2)[-3],
 \qquad
 H^1(C_{\overline X}^0\otimes^L\mathbb Z/2)=0.
\]
The sign-deleted choice $\mathbb Z$ has the same degree-one vanishing.
Hence (E.4.10)--(E.4.11) would imply
$H^1(\overline U_{\mathrm{et}},\mathbb Z/2)=0$.

There is an explicit nontrivial torsor contradicting that vanishing.
Put $R=\mathbb Z[1/2]$ and
\[
 B=R[t]/(t^2-2),\qquad \sigma(t)=-t.
\]
This is a finite free $R$-algebra of rank two, and it is \'etale
because $2t$ is a unit: $(2t)(t/4)=1$ in $B$.
The two factors $T-t$ and $T+t$ in $B[T]$ are comaximal, since
$2t$ is a unit. The Chinese remainder theorem therefore identifies
$B\otimes_R B$ with $B\times B$ by the two automorphisms
$1,\sigma$. Thus $\operatorname{Spec}B$ is a right torsor under the
constant group $\mathbb Z/2$ over $\operatorname{Spec}R$.
It has no section: a section would supply $a\in R\subset\mathbb Q$
with $a^2=2$, impossible because its $2$-adic valuation would satisfy
$2v_2(a)=1$. Equivalently $B$ has fraction field $\mathbb Q(\sqrt2)$
and is not the product algebra of a trivial torsor.

The extension is totally real and splits at the real completion:
$\mathbb Q(\sqrt2)\otimes\mathbb R\simeq\mathbb R\times\mathbb R$.
Attaching both real places therefore extends this torsor to the
Artin--Verdier site with the real point retained. It remains nontrivial,
since a section there would restrict to one over $R$.
Its class in $H^1(\overline U_{\mathrm{et}},\mathbb Z/2)$ is consequently
nonzero. This contradicts the required vanishing and proves the claim.
\end{proof}

The coefficient convention in (E.4.11) is concrete: in the proper
Flach--Morin theory, \citet[Proposition 6.11]{D_FlachMorin2018}
identifies $\mathbb Z(0)$ with the constant sheaf, and their
Lemma 3.13 identifies derived reduction with its finite-coefficient
\'etale cohomology. We have stated the analogous open comparison as
a requirement in order to prove its incompatibility with (E.4.10);
we have not misapplied their proper-scheme functoriality theorem to
the open immersion.

This counterexample explains a specific failed route toward the
all-twist integral extension. Geisser's Dedekind-ring purity theorem
has a truncation and a residue-characteristic qualification
\citep[Theorem 1.2(1)]{D_GeisserDedekind2004}. The missing qualification
cannot be removed in the stated coefficient theory. A corrected
integral localization theory must retain a nonzero contribution from
wild ramification at $2$, or alter (E.4.11), rather than use only the
finite-field complex $C_z^{-1}[-2]$ as its support term. Inverting $2$
removes this particular obstruction, and the rational result E.4g
is unaffected. The transcription 15:174 explicitly applies its displayed
purity sequence to $M=\mathbb Q(n)$, so this refutes an integral
extrapolation, not that rational assertion in the manuscript.

\paragraph{E.4k (Refuted): finite generation on opens together with the full coefficient comparison and rational restriction pattern.}
The obstruction persists without assuming the integral purity triangle.
There is no complex $C$ of abelian groups for which all three of the
following hold: $H^1(C)$ and $H^2(C)$ are finitely generated;
$H^1(C\otimes\mathbb Q)=H^2(C\otimes\mathbb Q)=0$;
and, for every $r\geq1$,
\[
 H^1(C\otimes^L\mathbb Z/2^r)
      \simeq H^1(\overline U_{\mathrm{et}},\mathbb Z/2^r),
 \qquad \overline U=
 \overline{\operatorname{Spec}\mathbb Z}-\{2\}.
 \tag{E.4.12}
\]
Thus the verified rational twist-zero restriction pattern cannot be
combined on this open with finite generation and every usual finite
coefficient comparison. No multiplicative or duality axiom is needed
for this refutation.

\begin{proof}
The first two assumptions make $A=H^1(C)$ and $B=H^2(C)$ finite.
The coefficient triangle supplies the exact sequence
\[
 0\longrightarrow A/2^rA
   \longrightarrow H^1(C\otimes^L\mathbb Z/2^r)
   \longrightarrow B[2^r]\longrightarrow0.
\]
Consequently the order of its middle group is at most $|A|\,|B|$,
independently of $r$.

We construct a class of exact order $2^r$ on the right of (E.4.12).
Set $m=r+2$ and let
\[
 K_m=\mathbb Q(\zeta_{2^m}+\zeta_{2^m}^{-1}),\qquad
 R=\mathbb Z[1/2],\qquad
 B_m=\bigl(R[T]/(T^{2^{m-1}}+1)\bigr)^{T\mapsto T^{-1}}.
\]
The polynomial is $\Phi_{2^m}(T)$; replacing $T$ by $T+1$ makes
it Eisenstein at $2$. The full cyclotomic algebra therefore has a
field as its generic fibre. It is finite \'etale over $R$, since
$T$ and $2^{m-1}T^{2^{m-1}-1}$ are units in the algebra.
The automorphisms $T\mapsto T^a$, for odd $a$ modulo $2^m$, give
its Galois torsor structure: on every geometric fibre the roots have
exact order $2^m$, and that group acts freely and transitively.
Taking the quotient by $a=-1$ shows that $\operatorname{Spec}B_m$
is a connected finite \'etale torsor under
\[
 G_m=(\mathbb Z/2^m\mathbb Z)^\times/\{\pm1\}
       \simeq\mathbb Z/2^{m-2}\mathbb Z.
\]
For the displayed cyclic identification one may use $5$ as a generator.
Indeed induction by squaring gives
$v_2(5^{2^s}-1)=s+2$ for every $s\geq0$, so $5$ has exact order
$2^{m-2}$ modulo $2^m$ and generates the units congruent to $1$
modulo $4$. Every odd residue is uniquely a sign times such a unit.

The generic field of $B_m$ is $K_m$, which is totally real.
Attaching all of its real places gives a connected torsor on
$\overline U_{\mathrm{et}}$ under $G_m\simeq\mathbb Z/2^r$.
Its monodromy character onto that cyclic group is surjective,
so its class in $H^1(\overline U_{\mathrm{et}},\mathbb Z/2^r)$
has exact order $2^r$. In particular the order of this cohomology
group is at least $2^r$, contradicting the preceding uniform bound
as $r$ increases. This proves the assertion.
\end{proof}

A possible integral extension retaining the rational pattern and all
these finite comparisons must therefore allow non-finitely-generated
torsion in the adjacent integral groups, or change one of those
requirements. For instance, a divisible $2$-primary torsion summand
can have nonzero $2^r$-torsion while having zero rationalization;
finite generation had excluded exactly this possibility.
The proposition does not assert nonexistence of all integral open
cohomology theories. It identifies the precise additional condition
that fails, beyond the already refuted exact proper sign table.


\paragraph{E.4l (Proved): the actual twist-zero support and every finite-prime open.}
The obstruction in E.4j has an explicit replacement. Throughout the next
five paragraphs put
\[
 X=\operatorname{Spec}\mathbb Z,\qquad
 U_S=\operatorname{Spec}\mathbb Z[S^{-1}],\qquad
 \overline U_S=\overline X-S,
\]
where $S$ is a finite set of rational primes; thus the real place is
retained. On $\overline U_S$ use the restriction of the Artin--Verdier
complex $\mathbb Z(n)_{\overline X}$ defined in
\citet[\S6.2 and Definition 6.5]{D_FlachMorin2018}, and write
\[
 A_S^n=R\Gamma(\overline U_{S,\mathrm{et}},
                     \mathbb Z(n)_{\overline X}|_{\overline U_S}).
\]
This notation specifies a complex, including its real-place convention.
It is not a claim that the manuscript defined these coefficients.
By Proposition 6.11 of that paper, $A_S^0$ is the cohomology of the
constant sheaf $\mathbb Z$. Define the discrete torsion group
\[
 D_p=\operatorname{Hom}_{\mathrm{cts}}
                   (\mathbb Z_p^\times,\mathbb Q/\mathbb Z).
\]
There is a local support isomorphism
\[
 R\Gamma_{\{p\}}(X_{\mathrm{et}},\mathbb Z)\simeq D_p[-3].
 \tag{E.4.13}
\]
Its tame subgroup is
$\operatorname{Hom}(\mathbb F_p^\times,\mathbb Q/\mathbb Z)$, induced by
reduction of units. The quotient is
$\operatorname{Hom}_{\mathrm{cts}}(1+p\mathbb Z_p,\mathbb Q/\mathbb Z)$.
For odd $p$ this quotient is $\mathbb Q_p/\mathbb Z_p$; for $p=2$ it is
$\mathbb Z/2\oplus\mathbb Q_2/\mathbb Z_2$.
Consequently, as abstract groups,
\[
 D_p\simeq
 \begin{cases}
  \mathbb Z/(p-1)\oplus\mathbb Q_p/\mathbb Z_p,&p\ne2,\\
  \mathbb Z/2\oplus\mathbb Q_2/\mathbb Z_2,&p=2.
 \end{cases}
 \tag{E.4.14}
\]
Only the finite tame summand is detected by
$C_p^{-1}[-2]$ from E.4a. The remaining summand is actual wild support.

For nonempty $S$, set
\[
 D_S^+=\operatorname{Hom}_{\mathrm{cts}}
 \left(
   \left(\prod_{p\in S}\mathbb Z_p^\times\right)/
                  \langle(-1)_{p\in S}\rangle,\ \mathbb Q/\mathbb Z
 \right).
\]
Then every cohomology group is given by
\[
 H^i(A_S^0)=
 \begin{cases}
  \mathbb Z,&i=0,\\
  D_S^+,&i=2,\\
  0,&\text{otherwise}.
 \end{cases}
 \qquad(S\ne\varnothing).
 \tag{E.4.15}
\]
For $S=\varnothing$, $A_\varnothing^0\simeq C_{\overline X}^0$
has the groups E.4b. In particular,
\[
 A_{\{2\}}^0\simeq
      \mathbb Z\oplus(\mathbb Q_2/\mathbb Z_2)[-2].
 \tag{E.4.16}
\]

\begin{proof}
Etale excision reduces (E.4.13) to the excellent henselian discrete
valuation ring $R=\mathbb Z_{(p)}^h$. The local duality pairing
of \citet[II, Theorem 1.8, with its pairing defined on p.~152]{D_MilneADT2006}
identifies $H^3_{\{p\}}(\operatorname{Spec}R,\mathbb Z)$ with the
continuous dual of the completion of $R^\times$, hence with $D_p$.
Here $\operatorname{Ext}^r_R(\mathbb Z,\mathbb G_m)
=H^r(R,\mathbb G_m)=0$ for $r>0$ by II, Proposition 1.5(a).
For completeness, the vanishing in all degrees and the shift can also
be read directly from localization. II, Corollary 1.2 gives
$H^0(R,\mathbb Z)=\mathbb Z$ and $H^2(R,\mathbb Z)=\mathbb Q/\mathbb Z$,
with the other groups zero. If $K=\operatorname{Frac}R$, local class
field theory and the sequence
$0\to\mathbb Z\to\mathbb Q\to\mathbb Q/\mathbb Z\to0$ give
\[
 H^2(K,\mathbb Z)
   =\operatorname{Hom}_{\mathrm{cts}}(\widehat{K^\times},
                                      \mathbb Q/\mathbb Z).
\]
The degree-two restriction from $R$ is the unramified character
subgroup, dual to the valuation $K^\times\to\mathbb Z$.
Also $H^3(K,\mathbb Z)=H^2(K,\mathbb Q/\mathbb Z)=0$:
finite local duality identifies the latter group's dual with the
inverse limit of finite groups of roots of unity in $K$, which is
zero because $K$ has only finitely many roots of unity.
Higher torsion cohomology vanishes by local cohomological dimension
two. The remaining quotient is precisely the dual of the unit
group. This proves (E.4.13), including that there are no omitted
cohomology groups. The equality of the henselian and completed
Galois groups is recalled in
\citet[I, proof of Theorem 2.14, pp.~36--37]{D_MilneADT2006}.

For odd $p$, the Teichmuller decomposition
$\mathbb Z_p^\times=\mu_{p-1}\times(1+p\mathbb Z_p)$ and the
$p$-adic logarithm identify the second factor with $\mathbb Z_p$.
For $p=2$, use
$\mathbb Z_2^\times=\{\pm1\}\times(1+4\mathbb Z_2)$ and the
logarithm on the second factor. This proves (E.4.14) and the
assertions about the tame subgroup.

The rational correction in Definition 3.6 of
\citet{D_FlachMorin2018} is zero in twist zero:
$R\Gamma(X,\mathbb Q(1))=0$, since $\mathbb Z(1)=\mathbb G_m[-1]$,
$\mathbb Z^\times$ and $\operatorname{Pic}(\mathbb Z)$ have zero
rationalization, and the higher $\mathbb G_m$ groups are torsion
by \citet[II, Proposition 2.1]{D_MilneADT2006}.
Thus $A_\varnothing^0=C_{\overline X}^0$.
Localization and E.4b now give, for any finite $S$,
\[
 0\longrightarrow H^2(A_S^0)\longrightarrow
       \bigoplus_{p\in S}D_p
   \xrightarrow{\partial_S}\mathbb Z/2
   \longrightarrow H^3(A_S^0)\longrightarrow0,
 \tag{E.4.17}
\]
as well as $H^0(A_S^0)=\mathbb Z$ and vanishing in all other degrees.

We identify the kernel and show surjectivity rather than assume
the boundary map. Constant rational coefficients have no positive
cohomology on $\overline U_S$: on the normal etale charts they are
the direct image of the constant rational sheaf at the generic
point, whose positive Galois cohomology vanishes by averaging over
finite quotients; the real stabilizer has the same vanishing.
It follows that
$H^2(A_S^0)=H^1(\overline U_{S,\mathrm{et}},\mathbb Q/\mathbb Z)$.
Finite abelian torsors on this Artin--Verdier site are the abelian
extensions of $\mathbb Q$ unramified outside $S$ that are totally
real. By the Kronecker--Weber theorem, with the ramification
control in its proof
\citep[I, Theorem 4.16 and Lemma 4.17, pp.~49--50]{D_MilneCFT2020},
their maximal Galois group is
\[
 \left(\prod_{p\in S}\mathbb Z_p^\times\right)/
                      \langle(-1)_{p\in S}\rangle.
\]
Indeed the maximal cyclotomic extension unramified outside $S$
has the product before quotienting, and complex conjugation is
the indicated diagonal element. Naturality of the coefficient
and localization sequences identifies the injection in (E.4.17)
with restriction of these characters to the local unit groups.
Its image is the annihilator of the diagonal element. The
quotient therefore has order two when $S\ne\varnothing$.
Consequently $\partial_S$ is surjective and is, under these
identifications, evaluation on that element. This proves
(E.4.15). For $S=\{2\}$ the quotient of the unit group by
$\{\pm1\}$ is $1+4\mathbb Z_2\simeq\mathbb Z_2$.
Finally the two nonzero cohomology groups in (E.4.16) split
uniquely in $D(\mathbb Z)$: the relevant ambiguity and obstruction
are $\operatorname{Ext}^2$ and $\operatorname{Ext}^3$, both zero.
\end{proof}

\paragraph{E.4m (Proved): the exact extra local term in every negative twist.}
For $n<0$, the coefficient convention of
\citet[\S3.1 and \S6.2]{D_FlachMorin2018} on $X_{\mathrm{et}}$ is
\[
 \mathbb Z(n)=
   \bigoplus_\ell j_{\ell,!}(\mathbb Q_\ell/\mathbb Z_\ell(n))[-1],
 \qquad j_\ell:X[1/\ell]\hookrightarrow X.
 \tag{E.4.18}
\]
With this convention there is a primary decomposition of the
actual support complex
\[
 R\Gamma_{\{p\}}(X_{\mathrm{et}},\mathbb Z(n))
 \simeq
 C_p^{\,n-1}[-2]\ \oplus\
 R\Gamma(\mathbb Q_p,\mathbb Q_p/\mathbb Z_p(n))[-2],
 \qquad n<0.
 \tag{E.4.19}
\]
In particular the second summand cannot be discarded when
preserving all finite coefficients.

\begin{proof}
Work again over $R=\mathbb Z_{(p)}^h$. For $\ell\ne p$ the
cyclotomic $\ell$-primary sheaf is unramified and the codimension-one
prime-to-$p$ purity calculation gives
\[
 R\Gamma_{\{p\}}(R,\mathbb Q_\ell/\mathbb Z_\ell(n)[-1])
  \simeq
 R\Gamma(\mathbb F_p,\mathbb Q_\ell/\mathbb Z_\ell(n-1))[-3].
\]
Equivalently, its local inertia computation uses the tame
$\mathbb Z_\ell(1)$ quotient: its first cohomology is the twist
$n-1$, its higher $\ell$-primary cohomology vanishes, and the
support and coefficient shifts supply the two additional degrees.
The sum of these terms is $C_p^{n-1}[-2]$ by the
negative-twist calculation E.4a. Cohomology and support commute
with these filtered sums; the local continuity statement is also
\citet[Lemma 2.1]{D_GeisserDedekind2004}.

For the $p$-primary summand put
$j:\operatorname{Spec}K\hookrightarrow\operatorname{Spec}R$.
The exact assertion
\[
 R\Gamma(R,j_!F)=0,\qquad
 R\Gamma_{\{p\}}(R,j_!F)\simeq R\Gamma(K,F)[-1]
\]
for an etale sheaf $F$ is
\citet[II, Proposition 1.1(a)]{D_MilneADT2006}.
Apply it to $F=\mathbb Q_p/\mathbb Z_p(n)$ and include the
shift $[-1]$ already present in (E.4.18). The result is
$R\Gamma(K,\mathbb Q_p/\mathbb Z_p(n))[-2]$.
The henselian-completion equivalence of finite extensions,
and then continuity for the torsion sheaf, replace $K$ by
$\mathbb Q_p$. This proves (E.4.19).
These computations take place at a finite closed point and
therefore are unchanged by the real-place correction.
\end{proof}

\paragraph{E.4n (Proved): full integral purity for positive twists at least two on these opens.}
The failure above does not extend to every weight. If $n\ge2$,
then for the closed immersion $i_p:\operatorname{Spec}\mathbb F_p
\hookrightarrow X$ the usual cycle-class purity morphism is a
quasi-isomorphism
\[
 \mathbb Z(n-1)_{\mathbb F_p}[-2]
       \xrightarrow{\ \sim\ } Ri_p^!\mathbb Z(n)_X.
 \tag{E.4.20}
\]
Consequently its derived global support is $C_p^{n-1}[-2]$,
including the characteristic-prime part, and restriction from
$\overline U_S$ to $\overline U_T$, $S\subset T$, has the
ordinary purity triangle in these weights.

\begin{proof}
There is a direct source with precisely the needed bound:
\citet[Proposition 5.1, p.~789]{D_GeisserDedekind2004}
proves full integral purity for $n\ge\operatorname{cd}_p(U_{\mathrm{et}})$,
where $U$ is the generic fibre of an essentially smooth scheme over
a discrete valuation ring of residue characteristic $p$.
Its historical Bloch--Kato hypothesis is now the norm-residue
theorem \citep[Theorems 6.16--6.17, pp.~436--437]{D_Voevodsky2011}.
By etale excision we may use the base ring
$R=\mathbb Z_{(p)}^h$ and the identity smooth morphism.
The generic field $K$ has the same absolute Galois group as
$\mathbb Q_p$, which has cohomological dimension two for every
torsion prime, including $p$. This follows from finite local
duality and is retained by finite extensions
\citep[I, Theorems 2.1 and 2.14]{D_MilneADT2006}.
Thus Geisser's inequality holds for every $n\ge2$.
His morphism is the ordinary purity morphism, so the resulting
local quasi-isomorphisms are the stalkwise isomorphisms in
(E.4.20), not a newly chosen map.

One can check the potentially delicate $p$-primary part separately.
\citet[Corollary 4.4, p.~786]{D_GeisserDedekind2004}
verifies the vanishing hypothesis in
\citet[Lemma 7.7, pp.~1543--1545]{D_FlachMorin2018}.
For this smooth discrete valuation ring its special fibre is a
normal crossing divisor. At each exponent $p^r$ the lemma's
localization proof gives
\[
 i_{p,*}\mathbb Z/p^r(n-1)[-2]\longrightarrow
 \mathbb Z/p^r(n)\longrightarrow
 \tau_{\le n}Rj_*\mu_{p^r}^{\otimes n}\longrightarrow.
 \tag{E.4.21}
\]
The argument is made with the cone of multiplication by $p^r$;
no inference from an isomorphism of pro-objects to each level
is required here. The first term vanishes because a perfect
field of dimension zero has
$W_r\Omega_{\mathbb F_p,\log}^{n-1}=0$ for $n-1>0$.
The stalks of $R^a j_*\mu_{p^r}^{\otimes n}$ for $a>1$
vanish: the strict henselian generic field has cohomological
dimension at most one
\citep[II, Lemma 1.3 and its proof, p.~150]{D_MilneADT2006}.
The truncation in (E.4.21) is therefore unnecessary.
It follows that $Ri_p^!\mathbb Z/p^r(n)=0$, exactly as
required by (E.4.20) at the residue-characteristic prime.
This also explains why the same argument does not give
twist-zero purity.

For $n-1>0$, ordinary etale motivic and Weil-etale cohomology
of the finite point agree. Indeed the comparison triangle of
\citet[Theorem 7.1]{D_Geisser2004} has the rational splitting
written in \citet[Proposition 3.17, equation (41), p.~1459]{D_FlachMorin2018}:
the rational Weil complex is the direct sum of the rational
etale complex and its shift by $[-1]$. E.4a makes that Weil
complex rationally zero, so both summands vanish. The integral
comparison triangle then gives the claimed isomorphism.
Hence the global support of (E.4.20) is the asserted
$C_p^{n-1}[-2]$.
\end{proof}

\paragraph{E.4o (Proved): the twist-one open complex has an actual finite-coefficient comparison.}
The complex constructed in E.4e can now be compared, on these
specific opens and with the specified real convention, to the
ordinary etale motivic complex. This is a sharper result than
merely comparing the displayed groups. For nonempty $S$ put
\[
 \begin{split}
 E_S&=\mathbb Z[S^{-1}]^\times=\{\pm1\}\oplus\mathbb Z^S,\\
 B_S&=\ker\left((\mathbb Q/\mathbb Z)^S
                     \xrightarrow{\sum}\mathbb Q/\mathbb Z\right),\\
 V_S&=\ker\left(\mathbb Q^S\xrightarrow{\sum}\mathbb Q\right),
 \qquad
 \Lambda_S=\ker\left(\mathbb Z^S\xrightarrow{\sum}\mathbb Z\right).
 \end{split}
 \tag{E.4.22}
\]
The unit identification uses the actual rational primes as
positive generators. There is an exact sequence
\[
 0\longrightarrow\Lambda_S\longrightarrow V_S
       \longrightarrow B_S\longrightarrow0.
 \tag{E.4.23}
\]
Then
\[
 A_S^1\simeq E_S[-1]\oplus B_S[-3],
 \qquad
 C_S^1\simeq
 \operatorname{Cone}\left(V_S[-3]\longrightarrow A_S^1\right),
 \quad S\ne\varnothing,
 \tag{E.4.24}
\]
where the arrow is zero on $E_S[-1]$ and the quotient in
(E.4.23) on $B_S[-3]$. The second isomorphism identifies this
cone with the complex in E.4e, compatibly with every restriction
$S\subset T$. In particular, for every positive integer $m$,
\[
 C_S^1\otimes^{\mathbb L}\mathbb Z/m
   \simeq
 R\Gamma(\overline U_{S,\mathrm{et}},\mathbb Z/m(1)),
 \tag{E.4.25}
\]
with $\mathbb Z/m(1)$ the derived reduction of the stipulated
Artin--Verdier complex.

\begin{proof}
Proposition 6.11 of \citet{D_FlachMorin2018} gives
$\mathbb Z(1)=\phi_*\mathbb G_m[-1]$ on the compactified site.
The global $\mathbb G_m$ calculation of
\citet[II, Proposition 2.1 and Remark 2.2(a),
pp.~163--165]{D_MilneADT2006} identifies the ordinary Brauer
group of $U_S$ with the kernel of the sum of invariants from
the removed finite places and the real place.
The real correction removes the real invariant. Explicitly,
the triangle from Proposition 6.10 of
\citet{D_FlachMorin2018} compares $\phi_*\mathbb G_m$ with
$R\phi_*\mathbb G_m$ and the positive-degree real cohomology.
The map $\operatorname{Br}(U_S)\to\operatorname{Br}(\mathbb R)
=\mathbb Z/2$ is surjective for $S\ne\varnothing$, since a
finite removed place can carry invariant $1/2$.
Its kernel is $B_S$. Higher groups cancel with their real
counterparts, while
$H^0(\overline U_S,\phi_*\mathbb G_m)=E_S$ and
$H^1(\overline U_S,\phi_*\mathbb G_m)
=\operatorname{Pic}(U_S)=0$.
Thus $A_S^1$ has only $E_S$ in degree one and $B_S$ in
degree three. Its splitting in (E.4.24) is unique because
$\operatorname{Ext}^{a}_{\mathbb Z}(-,-)=0$ for $a\ge2$.

To verify (E.4.23), lift an element of $B_S$ to a rational
vector. Its sum is integral; subtract that integer from one
coordinate to obtain a lift in $V_S$. The kernel is exactly
$\Lambda_S$. The cone in (E.4.24) consequently has $E_S$ in
degree one and $\Lambda_S$ in degree two, with no other
cohomology. Since $\Lambda_S$ is free, its splitting is also
unique: the possible ambiguity
$\operatorname{Ext}^1(\Lambda_S,E_S)$ is zero.
Its groups, identifications and restrictions are therefore
precisely those of the integral complex in E.4e.

For the proper member the corresponding data are
\[
 A_\varnothing^1\simeq
  (\mathbb Z/2)[-1]\oplus(\mathbb Q/\mathbb Z)[-4],
 \qquad
 \mathbb Q[-4]\longrightarrow A_\varnothing^1,
 \tag{E.4.26}
\]
where the map is the quotient $\mathbb Q\to\mathbb Q/\mathbb Z$
on the second summand. These are the Artin--Verdier trace and
the defining rational correction of
\citet[Theorem 3.5, Definition 3.6 and \S5.8.3]{D_FlachMorin2018}.
Its cone is $(\mathbb Z/2)[-1]\oplus\mathbb Z[-3]$.
Restriction from the proper member to a nonempty $S$ is the
unit inclusion and zero on the other displayed summands:
$\operatorname{Ext}^1(\mathbb Q/\mathbb Z,B_S)=0$ because
$B_S$ is divisible. Likewise the map of correction complexes
$\mathbb Q[-4]\to V_S[-3]$ is zero.
For two nonempty sets use the coordinate inclusions in
(E.4.22)--(E.4.23). These give commuting diagrams of actual
complexes and hence coherent cone maps, including the empty
set. They agree with E.4e; any possible off-diagonal map
between its displayed free cohomology groups vanishes.
Finally every correction complex here is a rational vector
space in one degree, so its derived reduction modulo $m$
is zero. Tensoring the defining cone triangle proves
(E.4.25), naturally for all these restriction maps.
\end{proof}

\paragraph{E.4p (Proved with an explicitly restricted scope): a coherent all-twist integral extension with the necessary wild terms.}
There is now an explicit family of bounded integral complexes,
with restrictions for all finite sets $S$ and every integer
twist, which has the corrected proper values, every finite
coefficient comparison, and exactly the rational values E.4g.
It is useful to state the construction and its limits precisely.
Let
\[
 D_X^n=R\operatorname{Hom}
        (R\Gamma(X,\mathbb Q(1-n)),\mathbb Q[-4])
\]
and let $\alpha_X^n:D_X^n\to A_\varnothing^n$ be the
proper correction of \citet[Theorem 3.5]{D_FlachMorin2018}.
For $n\ne1$ define
\[
 C_S^n=
 \operatorname{Cone}\left(
     D_X^n\xrightarrow{\alpha_X^n}A_\varnothing^n
                        \longrightarrow A_S^n\right).
 \tag{E.4.27}
\]
For $n=1$ use E.4o, including its proper member. Choose
functorial cochain representatives once, so all cones and
restriction maps in this definition are actual cochain cones.
For $S=\varnothing$ these are the proper Flach--Morin complexes.
For $n=0$ and $n\ge2$ the correction $D_X^n$ is zero;
thus $C_S^n=A_S^n$ in those weights.

For every $n,S$ and $m\ge1$ there are natural isomorphisms
\[
 C_S^n\otimes^{\mathbb L}\mathbb Z/m
       \simeq A_S^n\otimes^{\mathbb L}\mathbb Z/m
       \simeq
 R\Gamma(\overline U_{S,\mathrm{et}},\mathbb Z/m(n)).
 \tag{E.4.28}
\]
The first isomorphism follows because the correction is
uniquely divisible, and the second because derived reduction
is the cone of multiplication by $m$.
For a single additional removed prime $p$, localization has
the form
\[
 L_p^n\longrightarrow C_S^n\longrightarrow C_{S\cup\{p\}}^n
       \longrightarrow L_p^n[1],
 \qquad p\notin S,
 \tag{E.4.29}
\]
where
\[
 L_p^n\simeq
 \begin{cases}
 C_p^{n-1}[-2]\oplus
   R\Gamma(\mathbb Q_p,\mathbb Q_p/\mathbb Z_p(n))[-2],
       &n<0,\\
 D_p[-3],&n=0,\\
 C_p^0[-2],&n=1,\\
 C_p^{n-1}[-2],&n\ge2.
 \end{cases}
 \tag{E.4.30}
\]
For $n\ne1$ this is the actual etale localization triangle:
the fixed correction $D_X^n$ cancels in the morphism of
defining triangles. For $n=1$ it is E.4e and E.4o.
Multiple removals give the direct sum of these supports and
the compatible iterated localization triangles.

All objects in (E.4.27) are bounded. The proper objects and
their rational correction are bounded by the number-ring
calculation. The new local terms in (E.4.30) are bounded by
the local cohomological dimension bound and E.4a; localization
then proves the assertion for every finite $S$.
Finite generation is deliberately not an assertion:
(E.4.16) proves that it would be false.
Every support in (E.4.30) is rationally zero unless $n=1$.
Consequently (E.4.27) gives the entire rational table E.4g
and its restriction invariance, with the twist-one exception
already calculated in E.4c.

After inverting the residue characteristic $p$, the extra
summands in (E.4.30) vanish and the support becomes
$C_p^{n-1}[-2]\otimes\mathbb Z[1/p]$ in every twist.
In twist zero this follows from the exact unit-reduction
sequence in E.4l, whose quotient is $p$-primary; in negative
twists it follows directly from (E.4.19).
Thus the unrestricted integral purity statement is replaced
by an all-twist statement away from the residue
characteristic, and by the exact wild-support formula at it.

This construction settles existence and the finite-coefficient
comparison for this family of number-ring opens with the
stated coefficients and restrictions. It does not identify
the family with every proposed open Weil-etale topos, nor
construct an additional multiplicative theory on arbitrary
arithmetic schemes. The literal integral values requested in
Lead E2 remain refuted by E.4b and E.4d. The rational
manuscript localization is preserved, and is not the
integral purity assertion refuted in E.4j.


\subsection{Pro-fibre functors and structured crystalline fibres}
\label{sec:real-weil-E5}

\paragraph{E.5a (Known, with the required derivation): uniqueness over an algebraically closed field.}
Let $k$ be algebraically closed, and let $\mathcal C$ be a neutral
Tannakian category over $k$, without a tensor-generation or cardinality
restriction. Then any two $k$-valued fibre functors are isomorphic.
Consequently, whenever the pro-gerbe $T_\infty$ of 15:191 is an affine
Tannakian fibre gerbe and one such point is supplied by the indicated
realization, it has exactly one isomorphism class of $k$-points.

\begin{proof}
Given fibre functors $\omega,\eta$, the functor
$\operatorname{Isom}^{\otimes}(\omega,\eta)$ is a right fpqc torsor under
the affine group scheme $\operatorname{Aut}^{\otimes}(\omega)$, with
right action by precomposition. It represents coherent isomorphisms
on the whole category, not merely independent isomorphisms on its finite
subcategories. \citet[Theorem 1.1]{D_Wibmer2025} says that every such
torsor over an algebraically closed field has a $k$-point. That point is
exactly a tensor isomorphism $\omega\simeq\eta$. This is also Theorem 1.2
of the published paper (Corollary 1.2 of the arXiv version). Thus no
unjustified interchange of inverse limits and isomorphism classes is
needed.
\end{proof}

The existence of the realization is a separate hypothesis. In particular
a rational category with $\operatorname{End}(\mathbf1)=\mathbb Q$ cannot
have a $\mathbb Q$-linear fibre functor to $\operatorname{Vec}_k$ when
$\operatorname{char}k=p$: a unital map $\mathbb Q\to k$ is impossible.
The notation $T_\alpha(k)$ in characteristic $p$ requires an integral
model or a separately specified $k$-linear category, as suggested on
15:191. We do not obtain that model by reducing a rational fibre functor
formally.

\paragraph{E.5b (Refuted, with a perfect-field counterexample): passing to the limit does not remove $H^1$.}
Let
\[
 k=\bigcup_{m\geq0}\mathbb F_3(t^{1/3^m}),\qquad
 L=\bigcup_{m\geq0}\mathbb F_9(t^{1/3^m}),\qquad
 T=R^1_{L/k}\mathbb G_m.
\]
The fields are perfect and $L/k$ is a separable quadratic extension.
The right $T$-torsor $P$ defined by $N_{L/k}(z)=t$, with action by
multiplication, is nontrivial. The neutral category $\operatorname{Rep}_k(T)$
therefore has at least two nonisomorphic fibre functors, even for the
constant inverse system.

\begin{proof}
The $t$-adic valuation on $k$ has value group $\mathbb Z[1/3]$, and its
unique extension to this constant-field extension has the same value
group. The nontrivial automorphism of $L/k$ fixes that valuation. Hence
\[
 v(N_{L/k}z)=2v(z)\in2\mathbb Z[1/3].
\]
But $v(t)=1$ is not in $2\mathbb Z[1/3]$: an equality
$1=2a/3^m$ would make an odd integer even. Thus $P(k)$ is empty.
It becomes a trivial norm-one torsor after the separable extension $L$,
so it is a genuine smooth fpqc torsor. Twisting the forgetful functor
on $\operatorname{Rep}_k(T)$ by $P$ gives the second fibre functor;
an isomorphism to the forgetful functor would give a $k$-point of $P$.
Taking the same category at each index leaves this nonzero class intact.
\end{proof}

This counterexample refutes a universal perfect-field uniqueness
assertion for affine fibre gerbes. It is not asserted to construct
Grothendieck's unspecified integral motivic model over this field.
For a supplied neutral fibre functor the exact obstruction over an
arbitrary perfect field remains its ordinary fpqc $H^1(k,G)$ torsor
class; perfection alone gives no vanishing theorem.
This field has not been assumed to have cohomological dimension at
most one, so the example does not contradict the separate
Lang--Steinberg hypotheses discussed in R16.

\paragraph{E.5c (Proved and computed): different fibres can differ already as isocrystals.}
Over
$S=\operatorname{Spec}\mathbb F_5[a,(4a^3+27)^{-1}]$ consider the elliptic
scheme $y^2=x^3+ax+1$. Its fibres at $a=0,1$ are smooth. Direct enumeration
in the supplied program gives respectively $6$ and $9$ rational points.
Their Frobenius polynomials on $H^1$ are
\[
 T^2+5,\qquad T^2+3T+5.
 \tag{E.5}
\]
The first Newton polygon has slopes $1/2,1/2$; the second has slopes
$0,1$. Thus their rational contravariant crystalline realizations over
$\overline{\mathbb F}_5$ are nonisomorphic, and a fortiori their integral
$F$-crystals are nonisomorphic. After forgetting Frobenius, both underlying
$W(\overline{\mathbb F}_5)$-modules are free of rank two and are isomorphic.
The point counts, nonsingularity discriminants, and polygons give every
hypothesis of this example.

\paragraph{E.5d (Proved as a classification of the indicated structures): what the crystalline moduli mean.}
Let $k=\overline{\mathbb F}_p$, $W=W(k)$, $K_0=\operatorname{Frac}W$, and
let $\sigma$ be Witt Frobenius. For an unramified reductive group with
fixed smooth model $\mathcal G/\mathbb Z_p$, an isocrystal with a chosen
trivial underlying $G$-torsor is represented by $b\in G(K_0)$;
$F=b\sigma$. A change of trivialization $g$ changes $b$ to
$g^{-1}b\sigma(g)$. Isomorphism classes are therefore the corresponding
$\sigma$-conjugacy classes $B(G)$. This is the convention underlying
\citet[Introduction and \S\S3--4]{D_Kottwitz1985}.

If a dominant Hodge type $\mu$ is fixed, lattices in this fixed isocrystal
are the points
\[
 X_\mu(b)=\{g\mathcal G(W):
   g^{-1}b\sigma(g)\in\mathcal G(W)\mu(p)\mathcal G(W)\}.
 \tag{E.6}
\]
Here a coset represents the lattice $gW^d$ in a faithful representation
with its tensors. Applying $F$ to the lattice gives exactly the
double-coset condition displayed, so (E.6) follows directly by changing
a basis. This is the set defined in
\citet[\S1.2.9, author-PDF p.~14]{D_Kisin2017}.
Isomorphism classes after forgetting the embedding in the
fixed isocrystal require also quotienting on the left by
\[
 J_b(\mathbb Q_p)=\{h\in G(K_0):h^{-1}b\sigma(h)=b\};
\]
compare \citet[\S1.2.12, author-PDF p.~14]{D_Kisin2017}.
Thus a moduli problem with Frobenius and integral lattices is distinct
from $H^1$ of the ordinary group of an unstructured fibre functor.
Lead E3 captures one possible source of variation inside a Newton
stratum, but cannot explain all variation on 15:193: (E.5) varies the
Newton stratum itself. We have not identified the phrase
``la partie continue du $H^1$'' with a defined cohomology object in the
manuscript.

\paragraph{E.5e (Proved): unstructured integral fibre functors under a smooth model.}
Let $k$ be algebraically closed of characteristic $p$, let $W=W(k)$,
and let $\mathcal G/W$ be a smooth affine group scheme of finite
presentation. Then
$H^1_{\rm fpqc}(W,\mathcal G)=1$.
Indeed a $\mathcal G$-torsor $P$ is affine, smooth and of finite
presentation over $W$. Its special fibre is nonempty, because the
torsor map is faithfully flat, and has a $k$-point because $k$ is
algebraically closed. Smoothness lifts that point successively to
compatible points over $W/p^m$; the completeness of $W$ and finite
presentation of the affine scheme $P$ identify their inverse limit
with a $W$-point. A torsor with a point is trivial.

Consequently two existing $W$-valued fibre functors whose isomorphism
torsor has this group are isomorphic after forgetting $F,V$ and any
additional filtration. This proves precisely why ordinary smooth
torsor cohomology cannot explain nonisomorphism in that case. The
crystalline nonisomorphism in E.5c is a statement about $F$-crystals,
not a statement about those unstructured torsors. If the proposed
integral motivic group has no such smooth model, or if the
isomorphisms are required to commute with Frobenius, this vanishing
argument has different hypotheses and does not apply.

More generally, let $k$ be perfect of cohomological dimension at most
one, and assume in addition that the special fibre $\mathcal G_k$ is
connected. Then \citet[Theorem 1.9, p.~51]{D_Steinberg1965} gives
$H^1(k,\mathcal G_k)=1$. The same lifting argument proves
$H^1(W(k),\mathcal G)=1$. This applies in particular when $k$ is finite,
and also to a smooth affine group over $\mathbb Z_\ell$ with connected
special fibre. It verifies the finite-stage uniqueness assertion of
R16 with its actual hypotheses. It gives no such conclusion from
perfection alone, as E.5b shows. Even these finite-stage hypotheses do
not alone justify a passage to a pro-group, as the following example
shows.

\paragraph{E.5f (Refuted, including the $\operatorname{cd}(k)\leq1$ refinement): individually trivial finite restrictions need not admit compatible trivializations.}
Set
\[
 k=\bigcup_{m\geq0}\overline{\mathbb F}_3(t^{1/3^m}),\qquad
 G=\varprojlim_{n\geq0}(\mathbb G_{m,k},[2])
   =D_k(\mathbb Z[1/2]).
 \tag{E.5.1}
\]
This $k$ is perfect and has cohomological dimension at most one.
Indeed $\overline{\mathbb F}_3$ is algebraically closed and
$\operatorname{trdeg}_{\overline{\mathbb F}_3}k=1$, so the assertion
follows from \citet[Lemma 59.95.4, Tag 0F0T]{D_Stacks0F0T}.
Every finite stage $G_n=\mathbb G_m$ is smooth and connected and has
trivial $H^1(k,G_n)$.

For $n\geq0$ put $b_n=1$ if $3\mid n$ and $b_n=0$ otherwise. Take the
trivial right $G_n$-torsor $P_n=\mathbb G_{m,k}$, but give these torsors
the equivariant transition maps
\[
 P_{n+1}\longrightarrow P_n,\qquad
 x_{n+1}\longmapsto t^{b_n}x_{n+1}^2.
 \tag{E.5.2}
\]
Then $P=\varprojlim P_n$ is an actual affine fpqc $G$-torsor, and
$P(k)=\varnothing$. In particular the forgetful fibre functor on
$\operatorname{Rep}_k(G)$ and its twist by $P$ are isomorphic on every
finitely generated tensor subcategory but are not isomorphic on the
whole category.

\begin{proof}
The ring homomorphisms
$k[X_n,X_n^{-1}]\to k[X_{n+1},X_{n+1}^{-1}]$ send
$X_n$ to $t^{b_n}X_{n+1}^2$. They are injective and finite \'etale of
degree two. Their filtered union is nonzero, so its spectrum $P$ is
faithfully flat over the field $k$. The stagewise isomorphisms
\[
 P_n\times G_n\longrightarrow P_n\times P_n,\qquad
 (x,g)\longmapsto(x,xg)
\]
commute with the transition maps. Taking their affine inverse limits
gives $P\times G\simeq P\times P$. Thus $P$ satisfies the fpqc torsor
criterion, rather than merely specifying a system of unrelated
torsor classes. Extension of its structure group to $G_n$ gives $P_n$.

If $(x_n)$ were a point of $P(k)$, the $t$-adic valuations
$r_n=v_t(x_n)$ would belong to $\mathbb Z[1/3]$ and obey
$r_n=b_n+2r_{n+1}$. Hence, for every $N\geq1$,
\[
 r_0=\sum_{j=0}^{N-1}2^j b_j+2^Nr_N.
 \tag{E.5.3}
\]
The canonical embedding $\mathbb Z[1/3]\subset\mathbb Z_2$ makes
$2^Nr_N$ tend to zero: its $2$-adic valuation is at least $N$.
Therefore in $\mathbb Z_2$ we would have
\[
 r_0=\sum_{j\geq0}2^{3j}=\frac{1}{1-8}=-\frac17.
 \tag{E.5.4}
\]
But $-1/7\notin\mathbb Z[1/3]$, a contradiction.

A finite-dimensional representation of $D(\mathbb Z[1/2])$ has
finitely many weights, all in $2^{-n}\mathbb Z$ for some $n$. A finite
collection of such representations and all its tensor constructions
therefore factor through $G_n$. The two fibre functors restricted
there have isomorphism torsor $P_n$, which is trivial. An isomorphism
on the entire category, however, is a point of $P$, which we have
proved impossible.
\end{proof}

The obstruction is precisely a nonzero class of the sequence
$(t^{b_n})$ in
\[
 \begin{aligned}
 \varprojlim\nolimits^1(k^\times,[2])
 &=\operatorname{coker}(\Delta),\\
 \Delta:\prod_{n\geq0}k^\times&\longrightarrow\prod_{n\geq0}k^\times,
 & (x_n)&\longmapsto(x_nx_{n+1}^{-2}).
 \end{aligned}
 \tag{E.5.5}
\]
The displayed valuation proof proves nonzero directly, without
requiring a general cohomology comparison. The exact literature
analogue is \citet[Proposition 2.8 and Lemma 2.10, author-PDF
pp.~17--18]{D_Milne2003}: for countable inverse systems of tori with
surjective transition maps in characteristic zero, the kernel of the
map from fpqc $H^1$ to the
inverse limit of finite-stage $H^1$ is $\varprojlim^1 T_n(k)$.
The characteristic-zero hypothesis in that source is not silently
applied to (E.5.1); our characteristic-three proof is the preceding
explicit construction.

The same example lifts integrally. Over $W(k)$ replace $t^{b_n}$ in
(E.5.2) by the Teichm\"uller unit $[t]^{b_n}$. The coordinate transitions
are again finite \'etale and faithfully flat, and their limit is a
torsor under $D_{W(k)}(\mathbb Z[1/2])$. A $W(k)$-point would reduce to
a $k$-point of $P$, so the integral torsor is nontrivial, although all
its finite $\mathbb G_m$-torsor quotients are trivial.
This addresses the passage from E.5e to a pro-limit; it still makes
no assertion that this particular protorus is the unspecified
integral motivic group in the manuscript.

\paragraph{E.5g (Proved): conditions that do permit the passage to the limit.}
Let $(G_n,P_n)$ be a countable inverse system of affine groups and
equivariant torsors of finite type over a perfect field $k$ with
$\operatorname{cd}(k)\leq1$. Suppose each $G_n$ and every kernel of
$G_{n+1}\to G_n$ are smooth and connected, and the group transitions
are faithfully flat. Then each $P_n(k)$ is nonempty by
\citet[Theorem 1.9]{D_Steinberg1965}; the fibre over any point of
$P_n(k)$ is a torsor under the indicated kernel and so has a
$k$-point by the same theorem. Thus
$P_{n+1}(k)\to P_n(k)$ is surjective. Inductively choosing compatible
points proves $(\varprojlim P_n)(k)\ne\varnothing$.
The counterexample (E.5.1) has disconnected transition kernels
$\mu_2$, so it does not satisfy this sufficient condition.

There is also a compactness criterion, requiring no connectedness
condition on the transition kernels. For any directed inverse
system, if each $P_i(k)$ is a nonempty compact Hausdorff space and the
transition maps are continuous, the product of these spaces is
compact. The equations defining compatible points are closed and
have the finite intersection property: for finitely many equations
choose a common upper index and any point there. Their intersection
is nonempty. This applies to finite-type torsors over a finite
field once their finite-stage rational points exist, since these
point sets are finite. It also applies to affine finite-presentation
torsors over $\mathbb Z_\ell$ once their integral points exist:
each such point set is closed in some $\mathbb Z_\ell^N$ and hence
compact. Together with E.5e this proves pro-level uniqueness over
$\mathbb Z_\ell$ for an inverse presentation whose finite-stage
isomorphism torsors have smooth models with connected special fibres.
It supplies the missing compatibility argument in exactly those
cases, while (E.5f) proves that perfection and cohomological
dimension alone do not supply it.

\subsection{The local comparison class on Artin--Tate motives}
\label{sec:real-weil-E6}

\paragraph{E.6a (Proved): the class over $\mathbb Q_\ell$.}
Let $k$ be a number field and let $\mathcal{AT}(k)$ be the semisimple pure
$\mathbb Q$-linear category generated by finite Artin motives and
$\mathbb Q(1)$, not a category of mixed Tate extensions. Let $\Gamma_k$
be the absolute Galois group, regarded as the constant profinite group
scheme through its finite quotients. After fixing a geometric fibre
functor its Tannaka group is
\[
 G=\Gamma_k\times\mathbb G_m.
 \tag{E.7}
\]
For $k=\mathbb Q$, the right comparison torsor
$\operatorname{Isom}^{\otimes}(\omega_\ell,
\omega_{\mathrm{dR}}\otimes\mathbb Q_\ell)$ has class
\[
 [\iota_\ell,1]\in
 H^1(\mathbb Q_\ell,\Gamma_{\mathbb Q}\times\mathbb G_m),
 \tag{E.8}
\]
where $\iota_\ell:\Gamma_{\mathbb Q_\ell}\to\Gamma_{\mathbb Q}$ is a
local decomposition embedding, defined up to conjugacy. The convention
for (E.8) is that the descent cocycle of an Artin representation $V$ is
$\rho\circ\iota_\ell$. Reversing the orientation of the isomorphism
torsor reverses the cocycle convention; it does not change triviality.
Its Tate component is always trivial. Thus within this category a
nonzero class cannot occur independently of the Artin part.

\begin{proof}
Every pure object is a sum of Artin representations with an integer
Tate grading. Distinct gradings cannot be isomorphic because their
weights differ; the only common objects of the Artin and Tate
subcategories are sums of the unit. Tensor automorphisms are therefore
an automorphism on the finite Galois representations and an independent
scalar on $\mathbb Q(1)$, proving (E.7).
For an Artin representation $V$, de Rham realization is the invariant
space $(V\otimes\overline k)^{\Gamma_k}$, whereas the geometric \'etale
fibre is the underlying vector space. Base change to a separable closure
of a local field identifies them by Galois descent, and its descent
cocycle is the restriction of $\rho$ to the local decomposition group.
This proves the Artin component of (E.8) on every finite quotient,
and hence compatibly on the profinite group. A torsor under the scalar
factor is a one-dimensional vector space over the coefficient field;
choosing a nonzero vector trivializes it. Equivalently,
$H^1(\mathbb Q_\ell,\mathbb G_m)=0$ by Hilbert 90. This proves the last
assertion without choosing a canonical comparison of Tate bases.
\end{proof}

For example, on the quotient supplied by the Artin motive of
$\mathbb Q(i)/\mathbb Q$, the class at $\ell=3$ is the nontrivial
quadratic torsor $\operatorname{Spec}\mathbb Q_3(i)$: $-1$ reduces to
the nonsquare $2$ in $\mathbb F_3$. At $\ell=5$ it is trivial, because
$2^2\equiv-1\pmod5$ lifts by Hensel's lemma. This supplies a concrete
verification of the margin on 15:197 that the difference is already
visible on weight-zero Artin motives. Its degree is $H^1$, since it is
an isomorphism torsor; $H^2$ would classify the obstruction to existence
of the fibre gerbe itself.

\paragraph{E.6b (Proved): the number-field coefficient algebra.}
For general $k$,
$k\otimes_{\mathbb Q}\mathbb Q_\ell=\prod_{v\mid\ell}k_v$.
The comparison class is exactly the tuple
\[
 \bigl([\iota_v,1]\bigr)_{v\mid\ell},\qquad
 \iota_v:\Gamma_{k_v}\longrightarrow\Gamma_k,
 \tag{E.9}
\]
under the canonical product decomposition of $H^1$ over this finite
product of fields. A torsor over a disjoint union of spectra is a tuple
of torsors, which proves the asserted product decomposition. Applying
E.6a at each factor proves (E.9). If a rational point $x$ of a scheme
$S/k$ is used to specialize Artin local systems, replace $\iota_v$ by
its composite with the section $x_*:\Gamma_k\to\pi_1(S)$ and the finite
quotient defining that local system. This is the exact answer in terms
of the specified rational point; a numerical list of classes cannot be
given without specifying $S$, $x$ and its Artin quotient.
